\section{Setup}\label{sec:setup}
The data conventions preceding Theorem~\ref{thm:main} fix $k,D$,
the meromorphic domain, base point, logarithm branch and primitive
normalizations. They also define residues \eqref{eq:residue-vector},
the rate field \eqref{eq:exponential-field} and ranks \eqref{eq:rank-data}.
This section gives the coordinate dictionary used there.

\subsection{Admissible blocks}\label{subsec:blocks}
The five block types below have complex parameters $\lambda,\mu,c,s$,
with $c,s\ne0$.
\begin{description}\setlength{\itemsep}{0pt}\setlength{\parsep}{0pt}
\item[(B1)] $e^{\lambda z}$, $\lambda\in\C$.
\item[(B2)] A specified primitive of $\sum_{\mu\in M}q_\mu e^{\mu z}$,
$M$ finite, $q_\mu\in\C(z)$, with arbitrary complex additive constant.
The subclass $(\mathrm{B2}^E)$ consists of primitives with algebraic
rates, kernels $q_\mu\in\Qbar[z,z^{-1}]$, entire derivative and $F(0)=0$.
\par\noindent
The extension $(\mathrm{B2}^\tau)$ also permits finite towers. For each
$\eta\ne0$ and integer $h\ge1$, include levels $1\le k\le h$ defined by
\begin{equation}\label{eq:ein-tower}
\Ein_k(u)=\sum_{n\ge1}\frac{(-1)^{n+1}u^n}{n^kn!},\quad
D\Ein_k(\eta z)=\Ein_{k-1}(\eta z)/z\ (k\ge2).
\end{equation}
$D\Ein_1(\eta z)=(1-e^{-\eta z})/z$. Write $\Ein=\Ein_1$.
\item[(B3)] $\operatorname{Bess}(c,\nu)$, $c\ne0$, $\nu\in\Q$.
An integer Bessel order is \emph{resonant}. The local homogeneous equation has a logarithmic solution at zero. For order zero take $(I_c,V_c,S_c,W_c)$, where
\begin{equation}\label{eq:bessel-zero}
\begin{gathered}
I_c=\sum_{n\ge0}\frac{(cz/2)^{2n}}{(n!)^2},\quad
S_c=\sum_{n\ge0}\frac{\mathsf H_n(cz/2)^{2n}}{(n!)^2},\\
V_c=I_c',\ W_c=S_c',\ \mathsf H_n=\sum_{j=1}^n1/j,\ \mathsf H_0=0.
\end{gathered}
\end{equation}
As a function of its own argument, write $\mathsf S(u)=S_1(u)$; thus
$S_c(z)=\mathsf S(cz)$.
The operator $\mathscr B_{c,0}=zD^2+D-c^2z$ satisfies
$\mathscr B_{c,0}I_c=0$ and $\mathscr B_{c,0}S_c=2I_c'$. Thus $J_c=S_c-I_cL$
is a homogeneous solution. To define the other integer orders, put $R_n=a_n+b_nD$ for $n\ge0$, with
$(a_0,b_0)=(1,0)$, $(a_1,b_1)=(0,c^{-1})$ and
$(a_{n+1},b_{n+1})=(a_{n-1},b_{n-1})-2n(a_n,b_n)/(cz)$. Then
\begin{equation}\label{eq:bessel-integer}
I_{n,c}=R_nI_c=I_n(cz),\quad P_{n,c}=R_nS_c-b_nI_c/z,\quad
R_nJ_c=P_{n,c}-I_{n,c}L,
\end{equation}
with their derivatives, define the integer block. Use $I_{-n}=I_n$ and $P_{-n,c}=P_{n,c}$;
$P_{n,c}$ may have a pole at zero, e.g. $P_{1,c}=(W_c-I_c/z)/c$.
For every rational $\nu$ that is not a negative integer put
\begin{equation}\label{eq:bessel-nonresonant}
f_{\nu,c}=\sum_{n\ge0}\frac{(cz/2)^{2n}}{n!(\nu+1)_n},\qquad
\mathscr B_{c,\nu}=zD^2+(2\nu+1)D-c^2z.
\end{equation}
For every $\nu\in\Q\setminus\Z$, use the coordinate quadruple
$(f_{\nu,c},f'_{\nu,c},f_{-\nu,c},f'_{-\nu,c})$.
An order $\nu\notin\tfrac12\Z$ is \emph{nonresonant}: its two local
monodromy eigenvalues at zero are distinct.
The two solutions in this case are
$f_{\nu,c}=\Gamma(\nu+1)(cz/2)^{-\nu}I_\nu(cz)$ and
$z^{-2\nu}f_{-\nu,c}$. Half-integer orders reduce to exponentials.
The starting identities are
\[ f_{1/2,c}=\frac{\sinh(cz)}{cz},\qquad f_{-1/2,c}=\cosh(cz). \]
\item[(B4)] $\operatorname{Kum}(s,a)$, $s\ne0$, $a\notin\Z$, is resonant, with a
logarithmic homogeneous solution at zero. Its coordinates are
\[ (F,F',\Sigma,\Sigma')=(f_a(sz),D[f_a(sz)],S_a(sz),D[S_a(sz)]). \]
Here $f_a={}_1F_1(a;1;\cdot)$; the second Kummer parameter is fixed to one.
\begin{equation}\label{eq:kummer-block}
f_a(u)=\sum_{n\ge0}\frac{(a)_nu^n}{(n!)^2},\quad
S_a(u)=\sum_{n\ge0}\frac{(a)_nh_nu^n}{(n!)^2},\quad
h_n=2\mathsf H_n-\sum_{j<n}\frac1{j+a}.
\end{equation}
Here $(a)_0=1$ and $(a)_n=a(a+1)\cdots(a+n-1)$.
The operator $\mathscr K_{s,a}=zD^2+(1-sz)D-as$ satisfies
$\mathscr K_{s,a}F=0$ and $\mathscr K_{s,a}\Sigma=2F'-sF$. Thus $J=\Sigma-FL$ is
homogeneous. In the arithmetic statements we require $a\in\Q\setminus\Z$.
\item[(B5)] $L$, the primitive of $1/z$.
\end{description}
The tower at $\eta\ne0$ has depth $h_\eta$ in \eqref{eq:ein-tower};
closure is defined before Theorem~\ref{thm:main}.

The following are $E$-coordinates: listed exponentials with algebraic
rates, $(\mathrm{B2}^E)$ primitives, and tower levels $\Ein_k(\eta z)$
with $\eta\in\Qbar^\times$.
For Bessel blocks, require algebraic $c$ and rational order.
Include noninteger-order coordinates and the order-zero coordinates
$(I_c,V_c,S_c,W_c)$.
Other integer-order blocks contribute only the regular coordinates
$I_{n,c},I'_{n,c}$.
Kummer coordinates have algebraic $s$ and rational nonintegral $a$;
listed algebraic constants are also allowed. The integer-order
companions $P_{n,c},P'_{n,c}$ for $n\ge1$ are excluded and are handled
by order-zero reduction. The functions $L$, $G_{\eta,k}$ and
the logarithmic solutions $J$ are not $E$-coordinates.

Normalization, saturation, coordinate variables, duplicate/constant
graphs and Laurent elimination are defined together before
Theorem~\ref{thm:main}. Lemma~\ref{lem:exponential-log} proves
exponential independence.
